
\subsection*{The problem}

Large language models are increasingly deployed as autonomous agents that plan, call tools,
and keep long-running state.
Such runs fail often~\citep{cemri2025mast}, and they are hard to make durable.
A run may last minutes or hours and call external tools with real side effects, and when the
process crashes part-way through, the obvious recovery, re-running the agent from the start,
re-samples every model call.
The new run may take a different path, and it may repeat side effects that already happened.

Durable-execution engines, exemplified by Temporal~\citep{temporal2024docs}, solve the
engineering half of this problem by recording an event history.
Workflow code must be deterministic; calls to the outside world, including LLM calls, run as
\emph{Activities} whose results are recorded once and returned from the history on replay.
The idea is classical: logging every nondeterministic event so that a piecewise-deterministic
process can be replayed after a failure is the basis of message-logging recovery in
distributed systems~\citep{strom1985optimistic,elnozahy2002survey}.
For an LLM agent, however, this treats each model call as an opaque recorded value.
The engine makes replay correct but gives no account of the distribution the value was drawn
from, so it cannot say whether replay preserves the probability of a run, or what the
distribution of the rest of a run is after recovery.
Recent tools for LLM agents record and replay runs at their nondeterministic
boundaries~\citep{mudasiru2026agrepl,chawla2026chronicle} and check resume behaviour against
machine-checked contracts~\citep{khan2026resume,zheng2026execedits}, but they too are
deterministic or empirical.

Formal models of LLM agents, meanwhile, have focused elsewhere.
LLMbda~\citep{garby2026llmbda} gives a probabilistic, Lean-verified semantics for a single
agent's program with information-flow control, and $\lambda_A$~\citep{liu2026lambdaA} and
PPDL~\citep{mandel2026ppdl} give probabilistic semantics to agent programs and LLM flows.
The MSC calculus~\citep{bollig2026msc} proves coordination properties of many agents while
treating model outputs as arbitrary.
None of them models journaling, replay, or crash recovery.

\subsection*{Contribution}

We give a minimal probabilistic semantics for durable replay of an LLM agent and prove that
replay and crash recovery are sound in a precise probabilistic sense.

\begin{enumerate}

\item \textbf{A probabilistic semantics with a journal (Section~\ref{sec:semantics}).}
Every model call is a draw from an abstract probability measure and every response is recorded
in an append-only journal; the agent runs in record, replay, or resume mode.
Fixing the randomness on a tape lets runs in different modes be compared directly.

\item \textbf{Replay soundness (Theorem~\ref{thm:replay}).}
Replaying a run from its journal reproduces its trace and journal exactly.
Replay therefore succeeds with probability~1, and recording then replaying induces the same
distribution over traces as recording alone.

\item \textbf{Crash recovery (Corollary~\ref{cor:resume}).}
After a crash at any point, resuming with fresh randomness behaves exactly like an
uninterrupted run whose randomness is spliced at the crash point.
Resuming with the same randomness reproduces the original run, and resuming with fresh
randomness continues it with the original distribution conditioned on what survived, provided
the crash does not depend on a response that was never journaled.

\item \textbf{Many agents (Section~\ref{sec:multiagent}).}
Under any deterministic scheduler and with a single global journal, a system of concurrently
running agents that share memory and exchange messages unfolds into a single program tree, so
replay soundness and crash recovery carry over to whole systems (Theorem~\ref{thm:sys}).
We also prove a noninterference theorem for concurrently scheduled agents
(Theorem~\ref{thm:ni}): under a floating-label discipline, with labels on memory cells and
channels that never change, and a priority scheduler whose priorities depend only on public
events, the public trace of a system is the same whatever its secret data, pointwise for every
family of per-agent randomness tapes and hence in distribution.

\item \textbf{A Lean mechanization (Section~\ref{sec:mechanization}).}
The pointwise core of both results, and their generic lifting to any deterministic step
system, are proved in Lean~4 for every program tree, randomness tape, starting cursor, and
crash point, with no unproven steps; the instantiation to the multi-agent calculus and
the noninterference theorem are proved on paper.

\item \textbf{A prototype and measurements (Sections~\ref{sec:rust-runtime}
and~\ref{sec:evaluation}).}
A dependency-free Rust crate implements the semantics together with write-ahead logging for
tool side effects.
We measure replay determinism, recovery from a crash at every journal position, and journaling
overhead.

\end{enumerate}

\noindent
The paper's source, the Lean development, and the Rust prototype with all experiment
harnesses are available at \url{https://github.com/ankthba/agentical}.

\subsection*{Scope}

Logging nondeterministic events for replay is not new; what is new is a semantics in which
recovery with fresh randomness provably preserves the distribution of a run
(Corollary~\ref{cor:resume}(ii)), not only individual runs.
Replay never repeats a side effect, but whether recovery may perform an interrupted effect
again depends on the external world, which we do not formalize (Section~\ref{sec:effects}).
The multi-agent results assume deterministic scheduling and a single global journal; the
prototype runs a single agent. Section~\ref{sec:limitations} lists what lifting these
restrictions requires.

\subsection*{Roadmap}

Section~\ref{sec:background} reviews probabilistic semantics and durable execution.
Section~\ref{sec:semantics} defines the calculus, Section~\ref{sec:metatheory} proves replay
soundness and crash recovery, and Section~\ref{sec:multiagent} extends both to systems of
agents and proves noninterference.
Section~\ref{sec:rust-runtime} describes the prototype and Section~\ref{sec:evaluation} the
experiments.
Section~\ref{sec:related} discusses related work, Section~\ref{sec:limitations} limitations and
future work, and Section~\ref{sec:conclusion} concludes.
